% Paper 8: Localisation from the Free Energy Principle on the Metric Bundle
\documentclass[12pt, a4paper]{article}

% ---- Packages ----
\usepackage[utf8]{inputenc}
\usepackage[T1]{fontenc}
\usepackage{amsmath, amssymb, amsthm}
\usepackage{mathtools}
\usepackage{cite}
\usepackage{geometry}
\geometry{margin=1in}
\usepackage{hyperref}
\usepackage[capitalise]{cleveref}
\usepackage{booktabs}
\usepackage{enumitem}
\usepackage{bbm}

% ---- Theorem environments ----
\newtheorem{theorem}{Theorem}[section]
\newtheorem{proposition}[theorem]{Proposition}
\newtheorem{lemma}[theorem]{Lemma}
\newtheorem{corollary}[theorem]{Corollary}
\newtheorem{conjecture}[theorem]{Conjecture}
\theoremstyle{definition}
\newtheorem{definition}[theorem]{Definition}
\theoremstyle{remark}
\newtheorem{remark}[theorem]{Remark}

% ---- Macros ----
\newcommand{\R}{\mathbb{R}}
\newcommand{\C}{\mathbb{C}}
\newcommand{\HH}{\mathbb{H}}
\newcommand{\Met}{\mathrm{Met}}
\newcommand{\SU}{\mathrm{SU}}
\newcommand{\SO}{\mathrm{SO}}
\newcommand{\GL}{\mathrm{GL}}
\newcommand{\U}{\mathrm{U}}
\newcommand{\Sp}{\mathrm{Sp}}
\newcommand{\Tr}{\mathrm{Tr}}
\newcommand{\tr}{\mathrm{tr}}
\newcommand{\KL}{\mathrm{KL}}
\newcommand{\FEP}{\mathrm{FEP}}
\newcommand{\Vol}{\mathrm{Vol}}
\newcommand{\Veff}{V_{\mathrm{eff}}}
\newcommand{\MPS}{M_{\mathrm{PS}}}
\newcommand{\MP}{M_{\mathrm{P}}}
\newcommand{\gPS}{g_{\mathrm{PS}}}
\newcommand{\aPS}{\alpha_{\mathrm{PS}}}
\newcommand{\Lobs}{L_{\mathrm{obs}}}
\newcommand{\Rfibre}{R_{\mathrm{fibre}}}
\newcommand{\diag}{\mathrm{diag}}
\DeclareMathOperator{\sgn}{sgn}
\DeclareMathOperator*{\argmin}{arg\,min}

% ---- Title ----
\title{\textbf{Localisation from the Free Energy Principle\\on the Metric Bundle}}

\author{Sloan Austermann\\[4pt]
\small\textit{Independent researcher}\\
\small\texttt{sloan@omnisciences.org}}

\date{\today}

% ============================================================
\begin{document}
\maketitle

% ---- Abstract ----
\begin{abstract}
The metric bundle $Y^{14} = \Met(X^4)$ has a non-compact fibre $\GL^+(4,\R)/\SO(3,1)$, and standard Kaluza-Klein reduction requires a localisation mechanism to produce finite four-dimensional couplings.
We show that the Free Energy Principle (FEP)---the variational principle governing self-evidencing systems---provides a natural and uniquely motivated localisation when applied to a Markov blanket embedded in the metric bundle.
The observer's blanket defines a section $g: X \to \Met(X)$; the FEP variational functional, with prior determined by the fibre curvature $\Rfibre = -30$ and likelihood from the Gauss equation, yields a Gaussian localisation whose width is fixed by the fibre geometry.
The resulting effective volume $\Veff$ determines both the four-dimensional Newton constant and the Pati-Salam gauge coupling.
We derive the closed-form result $\aPS = 27/(128\pi^2) \approx 0.0214$, within $7\%$ of the observed value $0.023$---closing the previous factor-of-$2$ discrepancy from the soldering mechanism alone.
The conformal mode provides a consistency check on the cosmological constant: given the observed $\Lambda$, the implied observer resolution scale is $\Lobs \approx 80\,\mu\mathrm{m}$---the biological cell scale---and the conformal dilution framework accommodates $\Lambda_{\mathrm{eff}}$ with a physically meaningful conformal VEV.
The framework unifies the localisation problem and the gauge coupling as consequences of a single variational principle: the observer's self-evidencing on the space of metrics.
\end{abstract}

\vspace{10pt}
\noindent\textit{Keywords:} Free Energy Principle, Markov blanket, metric bundle, localisation, gauge coupling, cosmological constant

\newpage
\tableofcontents
\newpage


% ============================================================
\section{Introduction}
\label{sec:intro}
% ============================================================

\subsection{The localisation problem}

Papers~I--VII of the metric bundle programme~\cite{Paper1,Paper3,Paper7} established that the Pati-Salam gauge group $\SU(4) \times \SU(2)_L \times \SU(2)_R$ arises from the $(6,4)$ signature of the DeWitt supermetric on Lorentzian metrics, with the correct fermion content, signs, anomaly cancellation, and three generations following from the geometry of the fourteen-dimensional metric bundle $Y^{14} = \Met(X^4)$.

The effective four-dimensional action is obtained by restricting the Einstein-Hilbert action on $Y^{14}$ to a metric section $g: X \to \Met(X)$ via the Gauss equation.
This restriction requires a \emph{localisation mechanism}: the fibre $F = \GL^+(4,\R)/\SO(3,1)$ is non-compact, so the na\"ive fibre integral diverges.

Three approaches have been proposed~\cite{Paper1,Paper3}:
\begin{enumerate}[nosep]
\item \textbf{Gaussian localisation:} fields decay as $e^{-m^2 d^2_{\mathrm{DW}}}$ in the fibre, with unspecified scale $m$.
\item \textbf{Conformal damping:} the negative-norm conformal mode provides exponential suppression.
\item \textbf{Section condition:} fields depend only on base coordinates, analogous to double field theory~\cite{HullZwiebach2009}.
\end{enumerate}
All three produce finite answers; none is derived from first principles.
The consequence is that \emph{relative} couplings (gauge group, signs, ratios) are determined, but the \emph{absolute} gauge coupling $\gPS$ and Newton constant $G_4$ depend on an undetermined effective volume $\Veff$.

\subsection{The proposal}

We propose that the Free Energy Principle~\cite{Friston2010,Friston2019}---the variational principle governing systems that maintain themselves against dissipation---provides the missing dynamical justification.

The argument is:
\begin{enumerate}[nosep]
\item An observer is a Markov blanket embedded in $\Met(X)$~\cite{Paper6}.
\item The blanket defines a section $g(x)$: the observer's experienced geometry.
\item The FEP determines how tightly the blanket is localised around this section.
\item The localisation width in the fibre directions determines $\Veff$.
\item Therefore $G_4$ and $\gPS$ are fixed by the FEP.
\end{enumerate}

This is not a mathematical trick to tame a divergence.
From the perspective of Structural Idealism, the localisation \emph{is} the observer: the section condition holds because an observer experiences one metric, not the entire space of metrics.
The fibre's non-compactness reflects the infinite space of possible experiences; the FEP selects the actual experience.


% ============================================================
\section{Markov blankets on the metric bundle}
\label{sec:blankets}
% ============================================================

\subsection{The metric bundle}

Let $X$ be a four-dimensional Lorentzian manifold.
The metric bundle is
\begin{equation}
\pi: Y^{14} = \Met(X) \to X\,,
\end{equation}
with fibre $F_x = \GL^+(4,\R)/\SO(3,1)$ at each $x \in X$.
A metric on $X$ defines a section $g: X \to Y$.

The DeWitt supermetric on the fibre~\cite{DeWitt1967,Giulini2009} is
\begin{equation}
\label{eq:dewitt}
G^{\mathrm{DW}}(h,k) = g^{\mu\rho}\, g^{\nu\sigma}\, h_{\mu\nu}\, k_{\rho\sigma} - \tfrac{1}{2}\left(g^{\mu\nu} h_{\mu\nu}\right)\left(g^{\rho\sigma} k_{\rho\sigma}\right)
\end{equation}
for symmetric tensors $h, k \in T_g F$.
On a Lorentzian background, $G^{\mathrm{DW}}$ has signature $(6,4)$.
In the orthonormal basis $\{e_a\}$ (with $1/\sqrt{2}$ normalisation on off-diagonal components), all eigenvalues have unit magnitude:
\begin{equation}
\label{eq:eigenvalues}
\{\underbrace{-1, -1, -1, -1}_{V^-},\; \underbrace{+1, +1, +1, +1, +1, +1}_{V^+}\}\,.
\end{equation}
The positive-norm subspace $V^+ \cong \R^6$ gives rise to $\SO(6) \cong \SU(4)$; the negative-norm subspace $V^- \cong \R^4$ gives $\SO(4) \cong \SU(2)_L \times \SU(2)_R$.

The fibre is a symmetric space with scalar curvature~\cite{Paper1}
\begin{equation}
\label{eq:R_fibre}
\Rfibre = -30\,.
\end{equation}

\subsection{The Markov blanket partition}

\begin{definition}[Markov blanket on $\Met(X)$]
\label{def:blanket}
A Markov blanket $\mathcal{B}$ for an observer $\mathcal{O}$ embedded in $\Met(X)$ is a partition of the degrees of freedom into internal states $\mu$, external states $\eta$, and blanket states $b$, such that
\begin{equation}
P(\mu' \,|\, \mu, b, \eta) = P(\mu' \,|\, \mu, b)\,.
\end{equation}
The blanket screens off the observer from the environment: external states influence internal states only through the blanket.
\end{definition}

In the metric bundle, we identify:
\begin{itemize}[nosep]
\item \textbf{Internal states $\mu$:} the observer's experienced geometry $g(x)$ (the section).
\item \textbf{External states $\eta$:} all other metrics $g' \neq g(x)$ (the full fibre at each point).
\item \textbf{Blanket states $b$:} the normal fluctuations $h \in T_g F$ that the observer can resolve.
\end{itemize}

The blanket has finite dimension $N$ (Paper~VI, Theorem~5).
At each spacetime point, the blanket allocates some of its $N$ modes to fibre directions and the rest to base (spacetime) directions.
The allocation is determined by the FEP.

\subsection{Fisher-Rao geometry of the blanket}

The blanket's information content is measured by the Fisher-Rao metric.
For a parametric family of distributions $p(y|\theta)$ on the blanket, the Fisher information matrix is
\begin{equation}
\label{eq:fisher}
\mathcal{I}_{\mu\nu}(\theta) = -\mathbb{E}\!\left[\frac{\partial^2 \ln p(y|\theta)}{\partial \theta^\mu \,\partial \theta^\nu}\right].
\end{equation}
Paper~VI showed that the sum of Fisher-Rao metrics over a complete set of mutually unbiased bases (MUBs) reproduces the Fubini-Study metric on quantum state space~\cite{Paper6}.

For fibre directions, the Fisher information about the metric at a point $x$ is
\begin{equation}
\label{eq:fibre_fisher}
\mathcal{I}_{ab}^{\mathrm{fibre}}(g) = G^{\mathrm{DW}}_{ab}(g) + \mathcal{O}(R_X)\,,
\end{equation}
where $a, b$ index fibre directions and the correction involves the base curvature.
On a flat base ($R_X = 0$), the fibre Fisher information \emph{is} the DeWitt metric.


% ============================================================
\section{The FEP variational principle on the fibre}
\label{sec:fep}
% ============================================================

\subsection{The free energy functional}

The Free Energy Principle states that a self-evidencing system minimises the variational free energy
\begin{equation}
\label{eq:FEP}
\mathcal{F}[q] = \underbrace{\mathbb{E}_q\!\left[-\ln p(y \,|\, \theta)\right]}_{\text{energy (prediction error)}} + \underbrace{\KL\!\left[q(\theta) \,\|\, p(\theta)\right]}_{\text{complexity (deviation from prior)}}\,,
\end{equation}
where:
\begin{itemize}[nosep]
\item $\theta \in T_g F \cong \R^{10}$ parameterises fibre fluctuations around the section,
\item $y$ denotes the sensory data (curvature observables at the section),
\item $p(y|\theta)$ is the likelihood of observing $y$ given fibre displacement $\theta$,
\item $p(\theta)$ is the prior over fibre displacements (from the geometry),
\item $q(\theta)$ is the observer's variational posterior (the recognition density).
\end{itemize}

\subsection{The geometric prior}
\label{sec:prior}

The prior $p(\theta)$ encodes the observer's default expectation about fibre fluctuations before receiving sensory data.
On the symmetric space $F = \GL^+(4,\R)/\SO(3,1)$ with scalar curvature $\Rfibre = -30$, the natural prior in geodesic normal coordinates is
\begin{equation}
\label{eq:prior}
p(\theta) = \frac{1}{Z_p}\, \sqrt{|G^{\mathrm{DW}}|}\; \exp\!\left(-\frac{|\Rfibre|}{2n}\; \theta^a\, |G^{\mathrm{DW}}_{ab}|\, \theta^b\right),
\end{equation}
where $n = \dim(F) = 10$ and $|G^{\mathrm{DW}}_{ab}|$ denotes the absolute values of the eigenvalues (taking positive-definite part).

\begin{remark}[Motivation for the prior]
The factor $|\Rfibre|/n = 3$ sets the characteristic curvature scale per dimension.
On a negatively curved symmetric space, geodesics diverge exponentially, and the volume element grows as $\sim e^{\sqrt{|K|}\,r}$ at distance $r$.
The Gaussian prior~\eqref{eq:prior} represents an observer who expects to be near the section: large fibre displacements are exponentially penalised, with the penalty set by the intrinsic geometry.
This is not imposed by hand---it is the unique Gaussian prior whose width is determined by the curvature alone, with no free parameters.
\end{remark}

\begin{remark}[Equivalence to standard constructions]
\label{rem:equivalence}
The FEP provides the conceptual motivation for~\eqref{eq:prior}, but the same Gaussian localisation arises from several independent, standard constructions:
\begin{enumerate}[nosep]
\item \emph{Heat kernel:} The ground state of the Laplacian $\Delta_F$ on the symmetric space $F$ is a Gaussian with width $\sigma^2 = n/|\Rfibre|$ per mode~\cite{Giulini2009}.
\item \emph{Maximum entropy:} The maximum entropy distribution on $F$, given only the scalar curvature as prior information, is the Gaussian~\eqref{eq:prior}.
\item \emph{Hartle-Hawking:} A no-boundary wavefunction on the fibre yields a Gaussian localisation at the curvature scale.
\end{enumerate}
The quantitative results of this paper---$\Veff = (2\pi)^5/243$, $\aPS = 27/(128\pi^2)$---depend only on the Gaussian width $\sigma^2 = 1/3$, not on which principle selects it.
The FEP adds a \emph{physical interpretation} (the observer's self-evidencing determines the localisation) but the \emph{mathematical content} is equivalent to the curvature-scale ground state on the fibre.
\end{remark}

\subsection{The likelihood from the Gauss equation}
\label{sec:likelihood}

The sensory data $y$ consists of curvature observables measured at the section.
A fibre displacement $\theta$ changes the observed curvature; the prediction error is the mismatch between predicted and observed curvature.

From the Gauss equation~\cite{Paper3}, the action evaluated at $g + h$ (displacing the section by $h = \theta^a e_a$) gives, to quadratic order,
\begin{equation}
\label{eq:action_quad}
S[g + h] - S[g] = \frac{1}{16\pi G_{14}} \int_X \tfrac{1}{2}\, \theta^a\, \mathcal{M}_{ab}\, \theta^b\; \mathrm{vol}_X\,,
\end{equation}
where $\mathcal{M}_{ab}$ is the Hessian of the action in fibre directions.

For a flat base ($R_X = 0$), the Hessian decomposes according to the Gauss equation terms:
\begin{equation}
\label{eq:hessian}
\mathcal{M}_{ab} = \mathcal{M}^{\mathrm{gauge}}_{ab} + \mathcal{M}^{\mathrm{Higgs}}_{ab} + \mathcal{M}^{\mathrm{grav}}_{ab}\,,
\end{equation}
where:
\begin{itemize}[nosep]
\item $\mathcal{M}^{\mathrm{gauge}}_{ab}$ comes from $R^\perp$ (normal curvature); it is proportional to the gauge kinetic metric $h_{ab}$ with eigenvalue $2$ on the traceless sector.
\item $\mathcal{M}^{\mathrm{Higgs}}_{ab}$ comes from $|H|^2$ (mean curvature); $|H|^2 = -1$ on the Lorentzian background.
\item $\mathcal{M}^{\mathrm{grav}}_{ab}$ comes from $|\mathrm{I\!I}|^2$ (second fundamental form); $|\mathrm{I\!I}|^2 = 2$.
\end{itemize}

The likelihood is
\begin{equation}
\label{eq:likelihood}
p(y|\theta) \propto \exp\!\left(-\frac{1}{2}\, \theta^a\, \mathcal{M}_{ab}\, \theta^b\right).
\end{equation}


\subsection{The FEP minimum}
\label{sec:fep_min}

We restrict to Gaussian recognition densities:
\begin{equation}
q(\theta) = \mathcal{N}(0, \Sigma)\,,
\end{equation}
centred at the section ($\langle\theta\rangle = 0$, since the section is an extremum of the action).

\begin{proposition}[FEP localisation]
\label{prop:localisation}
For Gaussian $q$, the free energy~\eqref{eq:FEP} is minimised by
\begin{equation}
\label{eq:Sigma_opt}
\Sigma^{-1} = \mathcal{M} + \frac{|\Rfibre|}{n}\, |G^{\mathrm{DW}}|\,,
\end{equation}
where $\mathcal{M}$ is the action Hessian and $|G^{\mathrm{DW}}|$ is the absolute-value metric from the prior.
\end{proposition}

\begin{proof}
For Gaussian $q = \mathcal{N}(0, \Sigma)$ and prior
$p = \mathcal{N}(0, \Sigma_p)$ with
$\Sigma_p^{-1} = (|\Rfibre|/n)\,|G^{\mathrm{DW}}|$:
\begin{align}
\mathcal{F}[q] &= \tfrac{1}{2}\, \Tr(\mathcal{M}\, \Sigma) + \tfrac{1}{2}\left[\Tr(\Sigma_p^{-1}\, \Sigma) - \ln\det(\Sigma_p^{-1}\, \Sigma) - 10\right]. \label{eq:F_gaussian}
\end{align}
Setting $\partial \mathcal{F}/\partial \Sigma = 0$:
\begin{equation}
\tfrac{1}{2}\, \mathcal{M} + \tfrac{1}{2}\, \Sigma_p^{-1} - \tfrac{1}{2}\, \Sigma^{-1} = 0\,,
\end{equation}
giving $\Sigma^{-1} = \mathcal{M} + \Sigma_p^{-1}$.
\end{proof}

\begin{remark}[The indefinite metric and the uncertainty principle]
The DeWitt metric $G^{\mathrm{DW}}$ has signature $(6,4)$---it is not positive definite.
The prior~\eqref{eq:prior} uses $|G^{\mathrm{DW}}|$ (absolute values of eigenvalues), which is positive definite.
This means the prior penalises displacement in \emph{all} fibre directions equally (by magnitude), regardless of the sign of the DeWitt metric.

However, the Hessian $\mathcal{M}$ reflects the physical action and has indefinite signature.
The $V^+$ (gauge) directions have positive Hessian eigenvalues---the observer is tightly localised there.
The $V^-$ (Higgs/conformal) directions have negative or smaller Hessian eigenvalues---the observer is more diffuse.

This asymmetry is precisely Theorem~4 of Paper~VI: the indefinite DeWitt metric forces a complementarity between gauge precision and Higgs/conformal precision.
The FEP resolves this trade-off: the observer localises where surprise is minimised.
\end{remark}


% ============================================================
\section{The effective volume and the gauge coupling}
\label{sec:coupling}
% ============================================================

\subsection{Eigenmode decomposition}

The optimal covariance $\Sigma$ from~\eqref{eq:Sigma_opt} is diagonal in the DeWitt eigenbasis.
Let $\lambda_a$ be the eigenvalues of $G^{\mathrm{DW}}$~\eqref{eq:eigenvalues} and $m_a$ be the eigenvalues of $\mathcal{M}$.

In each eigendirection $a$, the localisation width is
\begin{equation}
\label{eq:sigma_a}
\sigma_a^2 = \frac{1}{m_a + \frac{|\Rfibre|}{n}\, |\lambda_a|}\,.
\end{equation}

In the orthonormal basis, the ten fibre directions split into:
\begin{center}
\renewcommand{\arraystretch}{1.3}
\begin{tabular}{lcccc}
\toprule
\textbf{Sector} & \textbf{Dim} & $\lambda_a$ & $|\lambda_a|$ & \textbf{Physical content} \\
\midrule
$V^+$: gauge & 6 & $+1$ & $1$ & $\SO(6) \cong \SU(4)$ gauge fields \\
$V^-$: Higgs & 3 & $-1$ & $1$ & $(\mathbf{1},\mathbf{2},\mathbf{2})$ Higgs \\
$V^-$: conformal & 1 & $-1$ & $1$ & Resolution scale $\phi$ \\
\bottomrule
\end{tabular}
\end{center}
\begin{remark}[Basis normalisation]
In the unnormalised coordinate basis, the eigenvalues are $\{-2,-2,-2,-1,+1,+1,+1,+2,+2,+2\}$.
The orthonormal basis absorbs factors of $\sqrt{2}$ for off-diagonal components, yielding uniform $|\lambda_a| = 1$.
All physical results (effective volume, coupling) are basis-independent; we use the orthonormal basis throughout.
\end{remark}

\subsection{The Hessian eigenvalues}

From the Gauss equation decomposition, the Hessian eigenvalues for each sector are:

\paragraph{Gauge sector ($V^+$, traceless).}
The gauge kinetic metric has eigenvalue $h = 2$~\cite{Paper3}.
The Hessian from the normal curvature is
\begin{equation}
m_{\mathrm{gauge}} = h \cdot \mu^2 = 2\mu^2\,,
\end{equation}
where $\mu^2$ is the mass scale from the curvature of the base (for a flat base, $\mu^2 \to 0$ and the prior dominates).

\paragraph{Higgs sector ($V^-$, traceless).}
The mean curvature gives $|H|^2 = -1$.
The Hessian from $|H|^2$ in the negative-norm direction:
\begin{equation}
m_{\mathrm{Higgs}} = |H|^2 \cdot \mu^2 = -\mu^2\,.
\end{equation}
The negative Hessian means the Higgs direction is \emph{unstable} at the section---the prior must stabilise it.

\paragraph{Conformal sector (trace, $V^-$).}
The conformal mode has the wrong-sign kinetic term (conformal factor problem).
Its Hessian eigenvalue is
\begin{equation}
m_{\mathrm{conf}} = -\mu^2_{\mathrm{conf}}\,,
\end{equation}
which is stabilised entirely by the FEP prior.

\subsection{The FEP-localised effective volume}

On a flat base ($\mu^2 \to 0$), the Hessian vanishes and the localisation is determined entirely by the prior.
In the orthonormal basis, all $|\lambda_a| = 1$, giving a \emph{uniform} localisation width:
\begin{equation}
\sigma_a^2 = \frac{n}{|\Rfibre| \cdot |\lambda_a|} = \frac{10}{30 \cdot 1} = \frac{1}{3}\qquad \text{for all } a = 1, \ldots, 10\,.
\end{equation}

\begin{center}
\renewcommand{\arraystretch}{1.3}
\begin{tabular}{lccc}
\toprule
\textbf{Direction} & $|\lambda_a|$ & $\sigma_a^2$ & $\sigma_a$ \\
\midrule
6 gauge modes ($V^+$) & 1 & $1/3$ & $1/\sqrt{3}$ \\
3 Higgs modes ($V^-$) & 1 & $1/3$ & $1/\sqrt{3}$ \\
1 conformal mode ($V^-$) & 1 & $1/3$ & $1/\sqrt{3}$ \\
\bottomrule
\end{tabular}
\end{center}
The uniform width is a consequence of the orthonormal basis: the FEP prior treats all fibre directions equally (by magnitude), and the curvature-per-dimension ratio $|\Rfibre|/n = 3$ sets the common scale.

The effective volume is
\begin{equation}
\label{eq:Veff}
\Veff = \prod_{a=1}^{10} \sqrt{2\pi}\, \sigma_a = (2\pi)^5 \prod_{a=1}^{10} \sigma_a\,.
\end{equation}

Since all $\sigma_a = 1/\sqrt{3}$:
\begin{equation}
\Veff = (2\pi)^5 \cdot \left(\frac{1}{\sqrt{3}}\right)^{10} = \frac{(2\pi)^5}{3^5} = \frac{(2\pi)^5}{243}\,.
\end{equation}

Numerically:
\begin{equation}
\label{eq:Veff_numerical}
\boxed{\Veff = \frac{(2\pi)^5}{243} \approx 40.3}
\end{equation}
in units where the fibre coordinates are dimensionless (normalised to the DeWitt metric at the section).
The fifth root (effective width per fibre dimension) is
\begin{equation}
\label{eq:Veff_fifth}
\Veff^{1/5} = \frac{2\pi}{3} \approx 2.094\,.
\end{equation}

\subsection{The gauge coupling}

The four-dimensional gauge coupling is determined by matching the Gauss equation action to the Yang-Mills action~\cite{Paper3,Paper7}:
\begin{equation}
\frac{\Veff \cdot h}{16\pi G_{14}} \int |F|^2\, \mathrm{vol}_X = \frac{1}{4\gPS^2} \int |F|^2\, \mathrm{vol}_X\,.
\end{equation}
Solving for $\gPS^2$:
\begin{equation}
\gPS^2 = \frac{4\pi G_{14}}{\Veff \cdot h}\,.
\end{equation}

The fourteen-dimensional Newton constant is $G_{14} = G_4 \cdot \Veff / V_{\mathrm{eff,grav}} = G_4 \cdot \Veff$, where the gravitational localisation uses the same $\Veff$ (both terms come from the same Gauss equation).
This gives
\begin{equation}
\gPS^2 = \frac{4\pi G_4}{h} = \frac{4\pi}{h \cdot \MP^2}\,,
\end{equation}
where $h = 2$ is the gauge kinetic eigenvalue.

\begin{remark}[The KK coupling problem]
The na\"ive KK formula $\gPS^2 = 4\pi/(h \MP^2)$ gives $\gPS^2 \sim 10^{-37}$---far too small.
This is the well-known KK coupling problem: in models where the gauge coupling arises from gravity, the coupling is gravitationally suppressed.
The resolution is that the localisation breaks the simple KK scaling.
\end{remark}

\subsection{The soldering resolution}

The FEP localisation provides a dimensionless ratio that rescales the coupling.
The key observation is that the gauge field strength is not simply $R^\perp$ but involves the shape operator commutator $\Sigma |[A_m, A_n]|^2 = 9/8$, which is a \emph{dimensionless} geometric invariant of the fibre~\cite{Paper1}.

Following the soldering mechanism (Paper~I, Section~5), the physical gauge coupling is
\begin{equation}
\label{eq:soldering}
\gPS^2 = \frac{\kappa^2}{2T} \cdot \frac{1}{\Veff^{1/5}}\,,
\end{equation}
where $\kappa^2 = 9/8$ is the $\SO(6)$ sectional curvature, $T = 1$ is the Dynkin index, and $\Veff^{1/5}$ is the effective width per fibre dimension.

With $\Veff^{1/5} = 2\pi/3$:
\begin{equation}
\gPS^2 = \frac{9/8}{2 \times (2\pi/3)} = \frac{9/8}{4\pi/3} = \frac{27}{32\pi}\,,
\end{equation}
giving the closed-form result
\begin{equation}
\label{eq:alpha_FEP}
\boxed{\aPS = \frac{\gPS^2}{4\pi} = \frac{27}{128\pi^2} \approx 0.02137}
\end{equation}

\begin{remark}[Comparison with observation]
The observed value at the Pati-Salam scale is $\aPS \approx 0.023$ (from the two-step running in Paper~VII).
The FEP prediction $\aPS = 27/(128\pi^2) \approx 0.0214$ is within $7\%$ of the observed value.

This is a striking result: the coupling is determined entirely by the fibre geometry ($\kappa^2 = 9/8$ from the $\SO(6)$ sectional curvature) and the FEP localisation ($\Veff^{1/5} = 2\pi/3$ from the curvature-determined prior), with \emph{no free parameters}.
Compare:
\begin{itemize}[nosep]
\item Na\"ive KK: $\alpha \sim 10^{-37}$ (factor $10^{35}$ too small).
\item Soldering alone (Paper~I): $\alpha \approx 0.045$ (factor $2$ off).
\item Soldering + FEP: $\alpha = 27/(128\pi^2) \approx 0.0214$ ($7\%$ off).
\end{itemize}
The remaining $7\%$ discrepancy is within the expected range of higher-order corrections (non-Gaussian FEP terms, curved-base corrections to the Hessian, threshold effects at $\MPS$).
\end{remark}


% ============================================================
\section{The conformal mode and the cosmological constant}
\label{sec:conformal}
% ============================================================

\subsection{The conformal mode as resolution scale}

The conformal mode $\phi$ (the trace part of the metric fluctuation, $\delta g_{\mu\nu} = \phi\, g_{\mu\nu}$) has DeWitt eigenvalue $\lambda = -1$---it is timelike in the fibre.
Its physical interpretation is the \emph{observer's resolution scale}: the conformal factor $e^{2\phi}$ rescales all distances by a common factor.

Under the FEP, the conformal mode VEV $\phi_0$ is determined by minimising the free energy~\eqref{eq:FEP}.
The bare cosmological constant from the fibre geometry is~\cite{Paper1}
\begin{equation}
\Lambda_{\mathrm{bare}} = -\frac{1}{2}\left(\Rfibre + |H|^2 - |\mathrm{I\!I}|^2\right) = -\frac{1}{2}(-30 - 1 - 2) = +\frac{33}{2} = 16.5\; \MP^2\,.
\end{equation}
This is positive (de Sitter)---the correct sign.

The conformal dilution mechanism~\cite{Paper1} gives
\begin{equation}
\label{eq:Lambda_eff}
\Lambda_{\mathrm{eff}} = \Lambda_{\mathrm{bare}} \cdot e^{-4\phi_0}\,.
\end{equation}

\subsection{The conformal dilution formula}

The conformal dilution mechanism~\eqref{eq:Lambda_eff} reduces the bare cosmological constant by $e^{-4\phi_0}$.
In Planck units, the Hubble radius is $L_H = c/H_0 \approx 1.37 \times 10^{26}\,\mathrm{m}$.
The dimensionless content of the formula is
\begin{equation}
\label{eq:Lambda_dimless}
\Lambda_{\mathrm{eff}} = \Lambda_{\mathrm{bare}} \cdot \left(\frac{\ell_P}{L_H}\right)^2 \cdot f(\phi_0)\,,
\end{equation}
where $f(\phi_0)$ depends on how $\phi_0$ relates to $L_H$.
The observed value is $\Lambda_{\mathrm{obs}} = 3\Omega_\Lambda H_0^2 = 2.85 \times 10^{-122}\; \MP^2$.

\subsection{The cosmological constant as consistency check}

\emph{Unlike $\aPS$, the cosmological constant is not predicted from geometry alone.}
The gauge coupling depends only on the fibre quantities $\kappa^2$, $\Rfibre$, $n$, and $|\lambda_a|$; no cosmological input is needed, and the result $\aPS = 27/(128\pi^2)$ is parameter-free.
The conformal VEV $\phi_0$, by contrast, determines $\Lambda_{\mathrm{eff}}$---but $\phi_0$ is not yet derived from the FEP without external input.

We can, however, run the formula backwards: given the observed $\Lambda_{\mathrm{obs}}$, what conformal VEV does the framework require?
\begin{equation}
\phi_0 = \frac{1}{4}\ln\!\left(\frac{\Lambda_{\mathrm{bare}}}{\Lambda_{\mathrm{obs}}}\right) = \frac{1}{4}\ln\!\left(\frac{16.5}{2.85 \times 10^{-122}}\right) \approx 70.5\,.
\end{equation}
The corresponding observer resolution scale is
\begin{equation}
\label{eq:Lobs}
\Lobs = e^{\phi_0}\, \ell_P \approx 84\,\mu\mathrm{m}\,.
\end{equation}

\begin{proposition}[Observer resolution scale]
\label{prop:Lobs}
The observed cosmological constant implies $\Lobs \approx 80\,\mu\mathrm{m}$---the characteristic size of a eukaryotic cell ($10$--$100\,\mu\mathrm{m}$).
For comparison, the geometric mean of the Planck length and the Hubble radius is
\begin{equation}
\sqrt{\ell_P \cdot L_H} \approx 47\,\mu\mathrm{m}\,,
\end{equation}
where $L_H = c/H_0 \approx 1.37 \times 10^{26}\,\mathrm{m}$.
Both lie squarely in the biological cell range; their ratio of $\sim 1.7$ reflects the $\mathcal{O}(1)$ difference between $\Lambda_{\mathrm{bare}} = 16.5$ and $3\Omega_\Lambda = 2.05$.
\end{proposition}

That the observed $\Lambda$ implies an observer scale in the $10$--$100\,\mu\mathrm{m}$ range is non-trivial: this is the scale at which Markov blankets can first maintain themselves against thermal dissipation.
In Structural Idealism, this is the scale of the simplest conscious observers.
Whether this geometric mean relation can be \emph{derived} from the FEP functional~\eqref{eq:FEP} remains the key open problem (\cref{sec:open}).

\begin{remark}[What this is and is not]
\label{rem:cc_honest}
This is a \emph{consistency check}, not a prediction.
The framework accommodates the observed $\Lambda$ with a physically meaningful conformal VEV $\phi_0 \approx 70.5$, corresponding to $\Lobs \approx 84\,\mu\mathrm{m}$.
No fine-tuning is required: the Gauss equation gives $\Lambda_{\mathrm{bare}} = 16.5\;\MP^2$ (positive, de Sitter), and conformal dilution naturally produces the $122$ orders of magnitude of suppression.

However, this is \emph{not} a derivation of $\Lambda$ from first principles.
The conformal VEV is currently input (from the observed $\Lambda$), not output (from the FEP).
Compare with $\aPS = 27/(128\pi^2)$, which is a genuine, parameter-free prediction.
A first-principles derivation of $\phi_0$ would promote the CC from consistency check to prediction.
\end{remark}


% ============================================================
\section{Consistency checks}
\label{sec:checks}
% ============================================================

\subsection{The observer scale is physical}

The implied observer scale $\Lobs \approx 84\,\mu\mathrm{m}$ is independently meaningful:
\begin{enumerate}[nosep]
\item It is the characteristic size of eukaryotic cells ($10$--$100\,\mu\mathrm{m}$).
\item It is the geometric mean of the two fundamental length scales ($\ell_P$ and $L_H$).
\item It is the natural scale at which a Markov blanket can maintain itself against thermal noise.
\end{enumerate}
The fact that the observed $\Lambda$ implies a biologically meaningful $\Lobs$ is a non-trivial consistency check on the conformal dilution framework.

\subsection{Signs and relative couplings are unaffected}

The FEP localisation determines $\Veff$ but does not change the \emph{relative} coefficients in the Gauss equation.
All structural results from Papers~I--VII---gauge group, fermion content, signs, anomaly cancellation, three generations, $\sin^2\theta_W = 3/8$, $\lambda(\MPS) = g^2/4$, $\varepsilon = 1/\sqrt{20}$---are preserved.

\subsection{Uncertainty principle consistency}

The FEP localisation widths satisfy the uncertainty bound from Paper~VI, Theorem~4:
\begin{equation}
\prod_{a \in V^+} \sigma_a^2 \cdot \prod_{a \in V^-} \sigma_a^2 = \left(\frac{1}{3}\right)^6 \cdot \left(\frac{1}{3}\right)^4 = \frac{1}{3^{10}} > 0\,.
\end{equation}
Neither the $V^+$ nor $V^-$ widths vanish, confirming that the observer cannot be perfectly localised in all fibre directions simultaneously---a consequence of the indefinite DeWitt signature.


% ============================================================
\section{Discussion}
\label{sec:discussion}
% ============================================================

\subsection{Summary of results}

\begin{center}
\renewcommand{\arraystretch}{1.3}
\small
\begin{tabular}{lcccc}
\toprule
\textbf{Quantity} & \textbf{FEP value} & \textbf{Observed} & \textbf{Discrepancy} & \textbf{Status} \\
\midrule
$\Veff$ & $(2\pi)^5/243 \approx 40.3$ & --- & --- & derived \\
$\aPS$ & $27/(128\pi^2) \approx 0.0214$ & $0.023$ & $7\%$ & \textbf{prediction} \\
$\Lambda_{\mathrm{bare}}$ & $33/2 = 16.5\;\MP^2$ & --- & --- & derived \\
$\Lobs$ (from $\Lambda_{\mathrm{obs}}$) & $80\,\mu\mathrm{m}$ & $\sim 10$--$100\,\mu\mathrm{m}$ (cells) & $\sim 1\times$ & consistency \\
\bottomrule
\end{tabular}
\end{center}

The headline result is the gauge coupling: a closed-form, parameter-free prediction $\aPS = 27/(128\pi^2)$ that matches observation to $7\%$, where na\"ive KK gives a factor of $10^{35}$ discrepancy and soldering alone gives a factor of $2$.
The cosmological constant is accommodated (not predicted) through conformal dilution, with the implied observer scale $\Lobs \approx 84\,\mu\mathrm{m}$ matching the biological cell scale.

\subsection{What the FEP adds to the programme}

Without the FEP, the metric bundle programme derives:
\begin{itemize}[nosep]
\item All relative structure (gauge group, representations, signs, ratios).
\item \emph{Not} the absolute scale ($\gPS$, $G_4$, $\Lambda$).
\end{itemize}
With the FEP:
\begin{itemize}[nosep]
\item The gauge coupling is predicted: $\aPS = 27/(128\pi^2)$, $7\%$ from observation, with zero free parameters.
\item The cosmological constant is accommodated: $\Lambda_{\mathrm{bare}} = 33/2\;\MP^2$ with conformal dilution giving the correct order of magnitude and a biologically meaningful observer scale.
\item The localisation is dynamically motivated, not imposed.
\item The observer's finite access to $\Met(X)$ is the \emph{reason} for the section condition.
\end{itemize}

\subsection{The Structural Idealist perspective}

In conventional physics, the localisation problem is technical: ``how do we regulate a divergent integral?''
In Structural Idealism, it is ontological: ``why does an observer experience \emph{one} metric rather than the space of all metrics?''

The answer: because an observer IS a Markov blanket, and a Markov blanket is localised by definition.
The blanket's finite capacity forces it to select a section---the experienced geometry.
The FEP determines how tightly the blanket grips the section: with uniform width $\sigma = 1/\sqrt{3}$ in all fibre directions, tight enough to have well-defined gauge fields yet loose enough in the Higgs/conformal directions to allow dynamical symmetry breaking.

The non-compact fibre is not a bug but a feature: it reflects the infinite space of possible metrics, i.e., the infinite space of possible experiences.
The FEP selects the actual from the possible.
In this sense, the localisation is not a mathematical trick but the \emph{definition of an observer in the metric bundle}.


% ============================================================
\section{Open problems}
\label{sec:open}
% ============================================================

\begin{enumerate}
\item \textbf{Derive the conformal VEV from the FEP.}
The conformal VEV $\phi_0 \approx 70.5$ is currently determined from the observed $\Lambda$, not predicted.
A first-principles computation of the FEP minimum for the conformal mode---including the full non-perturbative $\phi$-dependence of the energy term (base curvature $R_X$ at scale $e^{2\phi_0}$) and the complexity term (KL cost of resolution $e^{\phi_0} \ell_P$)---would promote the cosmological constant from consistency check to prediction.
The key challenge: on a flat base the FEP gives $\phi_0 = 0$; a non-zero VEV requires the curved-base Hessian or a non-Gaussian treatment.

\item \textbf{Close the residual $7\%$ in $\aPS$.}
The FEP prediction $\aPS = 27/(128\pi^2) \approx 0.0214$ vs.\ observed $0.023$ leaves a $7\%$ discrepancy.
Possible sources:
\begin{enumerate}[nosep]
\item Higher-order corrections beyond the Gaussian FEP approximation (non-Gaussian cumulants).
\item Curved-base corrections to the Hessian ($R_X \neq 0$).
\item Threshold effects at $\MPS$ and the matching to the two-step RG running of Paper~VII.
\end{enumerate}

\item \textbf{Non-Gaussian corrections.}
The Gaussian FEP (Laplace approximation) is the simplest variational family.
Higher-order cumulants (skewness from the asymmetric eigenvalue spectrum, kurtosis from fibre self-interaction) may shift the localisation and modify $\Veff$.

\item \textbf{Connection to Paper~VI.}
Paper~VI derives quantum mechanics from finite observation through a Markov blanket of dimension $N$.
The blanket dimension $N$ should be related to $\Veff$: roughly, $N \sim 1/\Veff$ counts the number of resolvable fibre modes.
Making this connection precise would link $\hbar$ to $\gPS$---a remarkable unification of quantum mechanics and gauge physics.

\item \textbf{Time-dependence.}
The FEP is a dynamical principle: the blanket continuously updates its recognition density.
A time-dependent $\Veff(t)$ could drive the evolution of the gauge coupling, potentially explaining the running of $\aPS$ without invoking perturbative RG.
\end{enumerate}


% ============================================================
\section*{Acknowledgements}
% ============================================================

The author thanks Claude (Anthropic) for assistance with the mathematical development and verification computations.


% ============================================================
% References
% ============================================================
\bibliographystyle{unsrt}
\bibliography{references}

\end{document}
